\section{总结与延伸}

\subsection{核心结论}

\begin{enumerate}
\item AI 社交的核心不是 AI 陪聊，而是让分身带着 context 主动创造真人连接。
\item Context 是 AI 社交的核心资产，智能可能平权，context 不平权。
\item 主体性是“做自己”的工程化表达，是分身能否代表用户行动的基础。
\item 隐私和 context 是交换关系，产品必须让用户在流失前感受到收益。
\item Elys 不是 AI 版 Tinder，目标是高维 context 匹配和 proactive 社交。
\item 模型公司未必天然赢得 AI 社交，应用公司的壁垒在 context、关系和产品心智。
\end{enumerate}

\subsection{开放问题}

\begin{itemize}
\item 用户是否愿意把足够多的 context 交给商业公司？
\item 分身替用户表达时，如何建立授权、撤回和责任机制？
\item AI 社交是否会增加真人连接，还是制造更多伪互动？
\item ChatGPT、Meta、豆包、Elys 这类玩家，谁更可能拥有长期 context 优势？
\item Proactive 产品如何避免变成骚扰和失控？
\end{itemize}

\subsection{拓展阅读}

\begin{itemize}
\item EP139 Agent 技术史：理解 proactive agent 的技术背景。
\item EP136 模型 OS：理解模型如何成为任务调度层。
\item 人机交互、社交网络、推荐系统与隐私计算相关资料：理解 AI 社交网络的多学科基础。
\item 电影 \textit{Her} 与 AI companion 产品讨论：理解陪伴叙事与现实产品的差异。
\end{itemize}

\begin{importantbox}{最后的判断}
AI 社交的真正难点不是让 AI 变得会聊天，而是让 AI 在尊重隐私和主体性的前提下，把一个人的 context 转化为更好的真人连接。如果它能做到，社交网络的底层单位会从“账号”变成“可行动的分身”。\end{importantbox}

\end{document}
